\section{从 CMU 到 OpenAI：职业选择}

\subsection{申请研究生：图灵奖推荐信也不保证 PhD}

翁家翌对 PhD 的申请并不顺利。虽然有 Yoshua Bengio 的推荐信（图灵奖得主），但暑研没有产出 paper，最终只拿到了 CMU 的 Master 而非 PhD。

\begin{knowledgebox}{清华的评价氛围}
在当时清华的氛围中，PhD 被认为优于 Master。翁家翌坦言自己"确实有一点失望"，花了一段时间调整。但他后来认为，\textbf{学历高低并不决定长期发展}，真正重要的是你的经验能不能匹配需求方的要求。"如果你想进工业界，那么读 PhD 就是浪费生命。"
\end{knowledgebox}

他在 CMU 期间在家上了一年网课（因为 COVID），但这段时间反而用来开发天授和 tuixue.online，把精力投入到更有长期价值的事情上。

翁家翌的申请经历提供了一个反直觉的案例：图灵奖得主的推荐信 + 清华计算机系本科 + OI 竞赛背景，依然可能只拿到 Master。他回忆说当时在知乎上看到很多人讨论申请结果，有人看到他"图灵奖强推"的背景却最后没有拿到 PhD，感到非常震惊——"竞争这么激烈"。

但翁家翌后来对此释然了。他认为 PhD vs Master 的高下之分是一种环境氛围制造的幻觉——"真正取决于你到底干什么"。他开始尝试挣脱这个评价体系，虽然在本科时期还没有完全挣脱出来，但已经意识到"应该创造自己的评价体系，而不是用其他人提供的评价体系"。

\subsection{CMU Master：疫情中的远程求学}

2020年秋季，翁家翌入学 CMU，但由于 COVID 疫情和签证问题，第一年完全在国内上网课。这段看似"浪费"的时间，他却充分利用——专注于天授的开发、tuixue.online 的维护，以及思考未来的职业方向。

他提到一个重要的心态转变：面对科伟（COVID）、国际政治动荡等宏大叙事，他选择\textbf{专注于手头上的事}——"不要天天去关注一些宏大的国际叙事，而是专注于手头上的事情，这样可能让自己内心更平静一些。"

\begin{knowledgebox}{疫情年的出国选择}
翁家翌那届（2019年12月申请）正好赶上 COVID。当时清华只有约 5\% 的人选择出国（比往年的 20\% 大幅下降），面临签证关闭、国际局势不确定等多重风险。翁家翌在这种不确定性下仍然坚持出国——与高中时选择清华降60分的决策模式如出一辙。
\end{knowledgebox}

\subsection{求职：DeepSeek vs OpenAI}

CMU 毕业时，翁家翌自嘲自己"开始时候吊儿郎当的"，投了18家公司，最后只收到了 Google 和 autoML（陈天琦老师的公司）的 offer。他不想去 Google——"在大厂当螺丝钉，做一些自己不是那么喜欢的事，比如前后端"。后来认真起来，又拿到了更多 offer。

他面临的核心选择是几家公司：

\begin{table}[H]
\centering
\begin{tabular}{lll}
\toprule
\textbf{公司} & \textbf{方向} & \textbf{结果} \\
\midrule
OpenAI & RL（John Schulman 组） & 接受 \\
幻方（后来的 DeepSeek） & RL Infra & 拿到 offer \\
NVIDIA & RL & 拿到 offer \\
Google & 通用 SWE & 拿到 offer \\
FAIR（Meta） & RL & 因流程原因被拒 \\
TikTok & -- & 拿到 offer \\
\bottomrule
\end{tabular}
\caption{翁家翌的求职选择}
\end{table}

\begin{importantbox}{选择 OpenAI 的逻辑}
这是 \textbf{before ChatGPT} 的时代。翁家翌选择 OpenAI 的原因是：（1）OpenAI 和 DeepMind 是当时 RL 领域最好的两个 research lab；（2）他想亲身体验世界最前沿的研究是怎么做的，而不是在学校"小作坊"里几个 PhD 手搓；（3）他想学习工业界有方法论地进行研究的方式。
\end{importantbox}

关于幻方（DeepSeek 前身），翁家翌透露当时幻方说要搞一个 AI Lab（即后来的 DeepSeek）。如果没有 OpenAI 的 offer，他会选择幻方。换言之，从"开天眼"的角度看，他当时面临的选择本质上是 \textbf{DeepSeek vs OpenAI}。

\subsection{John Schulman 的面试}

翁家翌被 John Schulman 亲自面试和招入 OpenAI。Schulman 认可他的核心原因是\textbf{GitHub 非常漂亮}——一个有良好工程能力的人对任何项目都有益。最后一轮面试是一道非常 end-to-end 的开放性题目，给了三个小时，翁家翌两个小时就完成了。

值得一提的是，这道面试题只测试过两个人——翁家翌和后来做 Codex 的同事（Andrew Y.），两人都通过了（"通过率百分之百"）。翁家翌对 Schulman 充满感激，甚至在 Schulman 离职那天难过了一个下午，关掉电脑什么都没做。

Schulman 赏识翁家翌的核心原因值得深思：不是论文，不是学历，而是\textbf{GitHub 非常漂亮}。这恰好验证了朱军老师提出的"GitHub 三位数 star"评价体系——在工业界，可见的工程产出确实比学术论文更有说服力。

\begin{knowledgebox}{翁家翌对读 PhD 的看法}
被问到找工作时是否考虑过读 PhD，翁家翌的回答很干脆：没有。"因为你接触了一些工业界的人会发现，如果你想进工业界那么读 PhD 就是浪费时间。你完全可以以 Master 为跳板，凑够 PhD 进工业界的标准。"他的策略是想清楚差异化——做出一些能让对方"挑选 Master 的你而不是另外一个 PhD"的项目。
\end{knowledgebox}

\subsection{本章小结}

翁家翌的求职经历体现了几个原则：（1）\textbf{差异化竞争}——用开源项目和工程能力弥补没有 PhD 的短板；（2）\textbf{面向需求优化}——搞清楚 AI Lab 需要什么样的人，然后让自己匹配；（3）\textbf{不被固有评价体系束缚}——Master 不如 PhD？那就在其他维度证明自己。


